\ifdefined\gkver\else\def\gkver{student}\fi
\documentclass[\gkver]{gaokaozhenti}
\begin{document}

\chapter{2004年全国理综Ⅱ}
%% paperType: 理综物理部分

\begin{enumerate}
\item
%% number: 14
%% paperName: 2004年全国理综Ⅱ第 14 题 6 分
%% typeId: 1
%% type: 单选题
%% chapter: 原子和原子核
%% point: 玻尔模型
%% method: 无
%% score: 6
%% degree: 650
%% duplicateId: 0
%% body:
现有 1200 个氢原子被激发到量子数为 4 的能级上，若这些受激氢原子最后都回到基态，则在此过程中发出的光子总数是多少？假定处在量子数为 $n$ 的激发态的氢原子跃迁到各较低能级的原子数都是处在该激发态能级上的原子总数的 $\frac{1}{n-1}$。\xzanswer{A}
\fourchoices[answer=A]
{2200}
{2000}
{1200}
{2400}

%% answer: A
%% memo:
\memoanswer{
	由题中假设可知，处于量子数 $n=4$ 的 1200 个氢原子跃迁到 $n=3$、$n=2$、$n=1$ 能级的原子数目均为 $1200\times\frac{1}{4-1}=400$ 个，此时发出光子 1200 个；处于量子数 $n=3$ 的 400 个氢原子跃迁到 $n=2$、$n=1$ 能级的原子数目均为 $400\times\frac{1}{3-1}=200$ 个，发出光子 400 个；处于量子数 $n=2$ 的原子总数为 $400+200=600$ 个，向 $n=1$ 跃迁发出光子 600 个，因此发出光子总数为 $1200+400+600=2200$ 个，故 A 正确。
}

\item
%% number: 15
%% paperName: 2004年全国理综Ⅱ第 15 题 6 分
%% typeId: 1
%% type: 单选题
%% chapter: 光学
%% point: 光的干涉
%% method: 无
%% score: 6
%% degree: 650
%% duplicateId: 0
%% body:
下面是四种与光有关的事实：①用光导纤维传播信号；②用透明的标准样板和单色光检查平面的平整度；③一束白光通过三棱镜形成彩色光带；④水面上的油膜呈现彩色。其中，与光的干涉有关的是\xzanswer{B}
\fourchoices[answer=B]
{①④}
{②④}
{①③}
{②③}

%% answer: B
%% memo:
\memoanswer{
	光导纤维利用光的全反射传播信号，白光通过三棱镜是光的折射引起的色散现象，检查光学平面的平整度和水面上的彩色油膜都是光的干涉现象，故 B 正确。
}

\item
%% number: 16
%% paperName: 2004年全国理综Ⅱ第 16 题 6 分
%% typeId: 1
%% type: 单选题
%% chapter: 气体性质
%% point: 热力学第一定律
%% method: 无
%% score: 6
%% degree: 650
%% duplicateId: 0
%% body:
一定量的气体吸收热量，体积膨胀并对外做功，则此过程的末态与初态相比\xzanswer{D}
\fourchoices[answer=D]
{气体内能一定增加}
{气体内能一定减小}
{气体内能一定不变}
{气体内能是增是减不能确定}

%% answer: D
%% memo:
\memoanswer{
	由热力学第一定律 $\Delta U=W+Q$ 知，气体吸热，$Q$ 为正值；对外做功，$W$ 为负值，但不知 $Q$ 和 $W$ 的数值大小，故内能的增减不能确定，故 D 正确。
}

\item
%% number: 17
%% paperName: 2004年全国理综Ⅱ第 17 题 6 分
%% typeId: 1
%% type: 单选题
%% chapter: 机械波
%% point: 波的多解性
%% method: 无
%% score: 6
%% degree: 550
%% duplicateId: 0
%% body:
如图，一简谐横波在 $x$ 轴上传播，轴上 a、b 两点相距 $12\Um$。$t=0$ 时 a 点为波峰，b 点为波谷；$t=0.5\Us$ 时，a 点为波谷，b 点为波峰。则下列判断正确的是\xzanswer{B}
\onepicture[w=7.0cm]{\includesvg{figs/17}}
\fourchoices[answer=B]
{波一定沿 $x$ 轴正方向传播}
{波长可能是 $8\Um$}
{周期可能是 $0.5\Us$}
{波速一定是 $24\Ums$}

%% answer: B
%% memo:
\memoanswer{
	因不知质点振动方向，无法确定波的传播方向，故 A 错误；由题意，$12=\left(n+\frac{1}{2}\right)\lambda\ (n=0,1,2,\ldots)$，$0.5=\left(k+\frac{1}{2}\right)T\ (k=0,1,2,\ldots)$，得 $\lambda=\frac{24}{2n+1}\Um\ (n=0,1,2,\ldots)$，$T=\frac{1}{2k+1}\Us\ (k=0,1,2,\ldots)$。当 $n=1$ 时，$\lambda=8\Um$，故 B 正确；周期 $T\neq0.5\Us$，故 C 错误；由 $v=\frac{\lambda}{T}=\frac{24(2k+1)}{2n+1}\Ums$ 知，只有当 $n=k$ 时 $v=24\Ums$，故 D 错误。
}

\item
%% number: 18
%% paperName: 2004年全国理综Ⅱ第 18 题 6 分
%% typeId: 1
%% type: 单选题
%% chapter: 力物体的平衡
%% point: 弹力
%% method: 无
%% score: 6
%% degree: 450
%% duplicateId: 0
%% body:
如图所示，四个完全相同的弹簧都处于水平位置，它们的右端受到大小皆为 $F$ 的拉力作用，而左端的情况各不相同：①中弹簧的左端固定在墙上，②中弹簧的左端受大小也为 $F$ 的拉力作用，③中弹簧的左端拴一小物块，物块在光滑的桌面上滑动，④中弹簧的左端拴一小物块，物块在有摩擦的桌面上滑动。若认为弹簧的质量都为零，以 $l_{1}$、$l_{2}$、$l_{3}$、$l_{4}$ 依次表示四个弹簧的伸长量，则有\xzanswer{D}
\onepicture[w=12.0cm]{\includesvg{figs/18}}
\fourchoices[answer=D]
{$l_{2}>l_{1}$}
{$l_{4}>l_{3}$}
{$l_{1}>l_{3}$}
{$l_{2}=l_{4}$}

%% answer: D
%% memo:
\memoanswer{
	由于弹簧质量为零，四种情况下，弹簧受力均为 $F$，故伸长量相同，$l_{1}=l_{2}=l_{3}=l_{4}$，故 D 正确。
}

\item
%% number: 19
%% paperName: 2004年全国理综Ⅱ第 19 题 6 分
%% typeId: 1
%% type: 单选题
%% chapter: 电磁感应
%% point: 切割磁感线电动势
%% method: 无
%% score: 6
%% degree: 550
%% duplicateId: 0
%% body:
一直升飞机停在南半球的地磁极上空。该处地磁场的方向竖直向上，磁感应强度为 $B$。直升飞机螺旋桨叶片的长度为 $l$，螺旋桨转动的频率为 $f$，顺着地磁场的方向看螺旋桨，螺旋桨按顺时针方向转动。螺旋桨叶片的近轴端为 a，远轴端为 b，如图所示。忽略 a 到转轴中心线的距离，用 $E$ 表示每个叶片中的感应电动势，则\xzanswer{A}
\onepicture[h=4.4cm]{\includesvg{figs/19}}
\fourchoices[answer=A]
{$E=\pi fl^{2}B$，且 a 点电势低于 b 点电势}
{$E=2\pi fl^{2}B$，且 a 点电势低于 b 点电势}
{$E=\pi fl^{2}B$，且 a 点电势高于 b 点电势}
{$E=2\pi fl^{2}B$，且 a 点电势高于 b 点电势}

%% answer: A
%% memo:
\memoanswer{
	电风扇叶片切割产生的感应电动势 $E=Bl\overline{v}=Bl\cdot\frac{\omega l}{2}$，又 $\omega=2\pi f$，解得 $E=\pi fl^{2}B$，由右手定则可判断 b 点电势比 a 点高，故 A 正确。
}

\item
%% number: 20
%% paperName: 2004年全国理综Ⅱ第 20 题 6 分
%% typeId: 1
%% type: 单选题
%% chapter: 电场
%% point: 电势能电势差电场力做功
%% method: 无
%% score: 6
%% degree: 550
%% duplicateId: 0
%% body:
如图一绝缘细杆的两端各固定着一个小球，两小球带有等量异号的电荷，处于匀强电场中，电场方向如图中箭头所示。开始时，细杆与电场方向垂直，即在图中Ⅰ所示的位置；接着使细杆绕其中心转过 $90^{\circ}$，到达图中Ⅱ所示的位置；最后使细杆移到图中Ⅲ所示的位置。以 $W_{1}$ 表示细杆由位置Ⅰ到位置Ⅱ过程中电场力对两小球所做的功，$W_{2}$ 表示细杆由位置Ⅱ到位置Ⅲ过程中电场力对两小球所做的功，则有\xzanswer{C}
\onepicture[h=4.8cm]{\includesvg{figs/20}}
\fourchoices[answer=C]
{$W_{1}=0$，$W_{2}\neq0$}
{$W_{1}=0$，$W_{2}=0$}
{$W_{1}\neq0$，$W_{2}=0$}
{$W_{1}\neq0$，$W_{2}\neq0$}

%% answer: C
%% memo:
\memoanswer{
	以两球整体为研究对象，杆由Ⅰ位置到Ⅱ位置，电场力对正、负带电小球都做正功，所以 $W_{1}\neq0$；杆由Ⅱ位置到Ⅲ位置，电场力对带正电小球做正功，对带负电小球做负功，且两功大小相等，所以 $W_{2}=0$，故 C 正确。
}

\item
%% number: 21
%% paperName: 2004年全国理综Ⅱ第 21 题 6 分
%% typeId: 1
%% type: 单选题
%% chapter: 运动和力
%% point: 牛顿定律中的图象
%% method: 无
%% score: 6
%% degree: 450
%% duplicateId: 0
%% body:
放在水平地面上的一物块，受到方向不变的水平推力 $F$ 的作用，$F$ 的大小与时间 $t$ 的关系和物块速度 $v$ 与时间 $t$ 的关系如图所示。取重力加速度 $g=10\Umsq$。由此两图线可以求得物块的质量 $m$ 和物块与地面之间的动摩擦因数 $\mu$ 分别为\xzanswer{A}
\twopicture[w=9.0cm]{figs/21a.png}{figs/21b.png}
\fourchoices[answer=A]
{$m=0.5\Ukg$，$\mu=0.4$}
{$m=1.5\Ukg$，$\mu=\frac{2}{15}$}
{$m=0.5\Ukg$，$\mu=0.2$}
{$m=1\Ukg$，$\mu=0.2$}

%% answer: A
%% memo:
\memoanswer{
	由图象可知 $0\sim2\Us$ 物体静止；$2\sim4\Us$，物体在外力 $F=3\UN$ 的作用下做匀加速运动，加速度为 $a=\frac{\Delta v}{\Delta t}=2\Umsq$，由牛顿第二定律得 $F-\mu mg=ma$，即 $3-10\mu m=2m$；$4\sim6\Us$，物体做匀速直线运动，故 $F'=\mu mg=2\UN$，即 $2=10\mu m$，联立解得 $m=0.5\Ukg$，$\mu=0.4$，故 A 正确。
}

\item
%% number: 22
%% paperName: 2004年全国理综Ⅱ第 22 题 18 分
%% typeId: 4
%% type: 实验题
%% chapter: 电路
%% point: 电路实验
%% method: 无
%% score: 18
%% degree: 550
%% duplicateId: 0
%% body:
用以下器材测量一待测电阻 $R_{x}$ 的阻值（$900\sim1000\UO$）：

电源 $E$，具有一定内阻，电动势约为 $9.0\UV$；

电压表 $\mathrm{V}_{1}$，量程为 $1.5\UV$，内阻 $r_{1}=750\UO$；

电压表 $\mathrm{V}_{2}$，量程为 $5\UV$，内阻 $r_{2}=2500\UO$；

滑线变阻器 $R$，最大阻值约为 $100\UO$；

单刀单掷开关 K，导线若干。
\begin{enumerate}
	\item
	测量中要求电压表的读数不小于其量程的 $\frac{1}{3}$，试画出测量电阻 $R_{x}$ 的一种实验电路原理图（原理图中的元件要用题图中相应的英文字母标注）。
	\item
	根据你所画的电路原理图在题给的实物图上画出连线。
	\item
	若电压表 $\mathrm{V}_{1}$ 的读数用 $U_{1}$ 表示，电压表 $\mathrm{V}_{2}$ 的读数用 $U_{2}$ 表示，则由已知量和测得量表示 $R_{x}$ 的公式为 $R_{x}=$ \_\_\_\_\_\_\_\_。
\end{enumerate}
\onepicture[align=right, w=6.0cm]{figs/22.png}

%% answer: 
\jdanswer{
\begin{enumerate}
	\item
	如图所示（电路原理图见详解）。

	\item
	如图所示（实物连线见详解）。

	\item
	$R_{x}=\frac{U_{1}r_{1}r_{2}}{U_{2}r_{1}-U_{1}r_{2}}$ 或 $R_{x}=\frac{U_{2}-U_{1}}{U_{1}}r_{1}$。
\end{enumerate}
}
%% memo:
\memoanswer{
	（1）在实验中测定的电阻 $R_{x}$ 的阻值（$900\sim1000\UO$）接近电压表 $\mathrm{V}_{1}$ 和 $\mathrm{V}_{2}$ 的内阻，属于测定大电阻，所以采用串联分压法，此外滑动变阻器 $R$ 的最大阻值很小，必须采用分压接法，故实验电路原理图如下方的图（a）或图（b）。图（a）中的电压表 $\mathrm{V}_{1}$ 和 $R_{x}$ 并联电阻约为 $420\UO$，两图都满足题中要求“电压表的读数不小于其量程的 $\frac{1}{3}$”。

	\onepicture[w=4.6cm]{figs/22_an1.png}

	或

	\onepicture[w=4.6cm]{figs/22_an2.png}

	（2）先连滑动变阻器的分压接法，然后再连接串联分压电路，如图所示。

	\onepicture[w=4.6cm]{figs/22_an3.png}

	或

	\onepicture[w=4.6cm]{figs/22_an4.png}

	（3）在图（a）中 $\frac{U_{1}}{r_{1}}+\frac{U_{1}}{R_{x}}=\frac{U_{2}}{r_{2}}$，在图（b）中 $U_{2}=U_{1}+R_{x}\times\frac{U_{1}}{r_{1}}$，化简得 $R_{x}=\frac{r_{1}r_{2}U_{1}}{r_{1}U_{2}-r_{2}U_{1}}$ 或 $R_{x}=\frac{(U_{2}-U_{1})r_{1}}{U_{1}}$。
}

\item
%% number: 23
%% paperName: 2004年全国理综Ⅱ第 23 题 16 分
%% typeId: 6
%% type: 计算题
%% chapter: 曲线运动
%% point: 平抛运动
%% method: 无
%% score: 16
%% degree: 450
%% duplicateId: 0
%% body:
一水平放置的水管，距地面高 $h=1.8\Um$，管内截面积 $S=2.0\Ucmq$。有水从管口以不变的速度 $v=2.0\Ums$ 源源不断地沿水平方向射出，设水流稳定后在空中有多少立方米的水。
%% answer: 
\jdanswer{
$V=2.4\times10^{-4}\Umc$
}
%% memo:
\memoanswer{
	以 $t$ 表示水由喷口处到落地所用的时间，有

	\[ h=\frac{1}{2}gt^{2} \tag{①} \]

	单位时间内喷出的水量为

	\[ Q=Sv \tag{②} \]

	空中水的总量应为

	\[ V=Qt \tag{③} \]

	由以上各式得

	\[ V=Sv\sqrt{\frac{2h}{g}} \tag{④} \]

	代入数值得 $V=2.4\times10^{-4}\Umc$ ⑤。
}

\item
%% number: 24
%% paperName: 2004年全国理综Ⅱ第 24 题 18 分
%% typeId: 6
%% type: 计算题
%% chapter: 磁场
%% point: 带电粒子在组合场中的运动
%% method: 无
%% score: 18
%% degree: 550
%% duplicateId: 0
%% body:
如图所示，在 $y>0$ 的空间中存在匀强电场，场强沿 $y$ 轴负方向；在 $y<0$ 的空间中，存在匀强磁场，磁场方向垂直 $xy$ 平面（纸面）向外。一电量为 $q$、质量为 $m$ 的带正电的运动粒子，经过 $y$ 轴上 $y=h$ 处的点 $P_{1}$ 时速率为 $v_{0}$，方向沿 $x$ 轴正方向；然后，经过 $x$ 轴上 $x=2h$ 处的 $P_{2}$ 点进入磁场，并经过 $y$ 轴上 $y=-2h$ 处的 $P_{3}$ 点。不计重力。求：
\begin{enumerate}
	\item
	电场强度的大小。
	\item
	粒子到达 $P_{2}$ 时速度的大小和方向。
	\item
	磁感应强度的大小。
\end{enumerate}
\onepicture[align=right, h=6.0cm]{figs/24.png}
%% answer: 
\jdanswer{
\begin{enumerate}
	\item
	$E=\frac{mv_{0}^{2}}{2qh}$。

	\item
	$v=\sqrt{2}v_{0}$，速度方向与 $x$ 轴夹角为 $45^{\circ}$。

	\item
	$B=\frac{mv_{0}}{qh}$。
\end{enumerate}
}
%% memo:
\memoanswer{
	（1）粒子在电场、磁场中运动的轨迹如图所示。设粒子从 $P_{1}$ 到 $P_{2}$ 的时间为 $t$，电场强度的大小为 $E$，粒子在电场中的加速度为 $a$，由牛顿第二定律及运动学公式得

	\[ qE=ma \tag{①} \]

	\[ v_{0}t=2h \tag{②} \]

	\[ \frac{1}{2}at^{2}=h \tag{③} \]

	\onepicture[w=4.6cm]{figs/24_an.png}

	联立解得

	\[ E=\frac{mv_{0}^{2}}{2qh} \tag{④} \]

	（2）粒子到达 $P_{2}$ 时速度沿 $x$ 方向的分量仍为 $v_{0}$，以 $v_{1}$ 表示速度沿 $y$ 方向分量的大小，$v$ 表示速度的大小，$\theta$ 表示速度和 $x$ 轴的夹角，则有

	\[ v_{1}^{2}=2ah \tag{⑤} \]

	\[ v=\sqrt{v_{1}^{2}+v_{0}^{2}} \tag{⑥} \]

	\[ \tan\theta=\frac{v_{1}}{v_{0}} \tag{⑦} \]

	由②③⑤式得 $v_{1}=v_{0}$⑧，

	由⑥⑦⑧式得 $v=\sqrt{2}v_{0}$⑨，

	\[ \theta=45^{\circ} \tag{⑩} \]

	（3）设磁场的磁感应强度为 $B$，在洛伦兹力作用下粒子做匀速圆周运动，由牛顿第二定律得

	\[ qvB=m\frac{v^{2}}{r} \tag{\textcircled{\scriptsize 11}} \]

	$r$ 是圆周的半径。此圆周与 $x$ 轴和 $y$ 轴的交点分别为 $P_{2}$、$P_{3}$，因 $OP_{2}=OP_{3}$，$\theta=45^{\circ}$，由几何关系可知，连线 $P_{2}P_{3}$ 为圆轨道的直径，由此可求得

	\[ r=\sqrt{2}h \tag{\textcircled{\scriptsize 12}} \]

	由⑨\textcircled{\scriptsize 11}\textcircled{\scriptsize 12}可得

	\[ B=\frac{mv_{0}}{qh} \tag{\textcircled{\scriptsize 13}} \]
}

\item
%% number: 25
%% paperName: 2004年全国理综Ⅱ第 25 题 20 分
%% typeId: 6
%% type: 计算题
%% chapter: 运动和力
%% point: 动量和能量
%% method: 无
%% score: 20
%% degree: 350
%% duplicateId: 0
%% body:
柴油打桩机的重锤由气缸、活塞等若干部件组成，气缸与活塞间有柴油与空气的混合物。在重锤与桩碰撞的过程中，通过压缩使混合物燃烧，产生高温高压气体，从而使桩向下运动，锤向上运动。现把柴油打桩机和打桩过程简化如下：

柴油打桩机重锤的质量为 $m$，锤在桩帽以上高度为 $h$ 处（如图\ref{2004全国理综Ⅱ25a}）从静止开始沿竖直轨道自由落下，打在质量为 $M$（包括桩帽）的钢筋混凝土桩子上。同时，柴油燃烧，产生猛烈的推力，锤和桩分离，这一过程的时间极短。随后，桩在泥土中向下移动一距离 $l$。已知锤反跳后到达最高点时，锤与已停下的桩帽之间的距离也为 $h$（如图\ref{2004全国理综Ⅱ25b}）。已知 $m=1.0\times10^{3}\Ukg$，$M=2.0\times10^{3}\Ukg$，$h=2.0\Um$，$l=0.20\Um$，重力加速度 $g=10\Umsq$，混合物的质量不计。设桩向下移动的过程泥土对桩的作用力 $F$ 是恒力，求此力的大小。
\twopicture[w=5.2cm, labelA=2004全国理综Ⅱ25a, labelB=2004全国理综Ⅱ25b]{figs/25a.png}{figs/25b.png}
%% answer: 
\jdanswer{
$F=2.1\times10^{5}\UN$
}
%% memo:
\memoanswer{
	重锤做自由落体运动，碰桩前的速度大小为 $v_{1}$，方向竖直向下，由运动学公式得

	\[ v_{1}=\sqrt{2gh} \tag{①} \]

	设碰后瞬间重锤的速度大小为 $v_{2}$，已知锤上升高度为 $(h-l)$，由运动学公式得

	\[ v_{2}=\sqrt{2g(h-l)} \tag{②} \]

	设碰后桩的速度大小为 $v$，方向竖直向下，由动量守恒得

	\[ mv_{1}=Mv-mv_{2} \tag{③} \]

	桩下降的过程中，由功能关系得

	\[ \frac{1}{2}Mv^{2}+Mgl=Fl \tag{④} \]

	由①②③④式得

	\[ F=Mg+\frac{mg}{l}\cdot\frac{m}{M}\left[2h-l+2\sqrt{h(h-l)}\right] \]

	代入数值得 $F=2.1\times10^{5}\UN$。
}

\end{enumerate}

\end{document}
